\chapter{Green Functions for PDE}
\label{ch-green-fun-pde}

\section{Introduction}
Let
\beq
\call = \partial_t - D \partial^2_x
\eeq
The Diffusion Equation (DE) in one space dimension $x$,
 with a source $S(x,t)$, is
\beq
\call \eta(x,t) = S(x,t)
\eeq
We will illustrate the use of Green
functions by solving this PDE
for all $x>0$ and $t>0$, with boundary conditions (BC)
at $x=0$.

Common types of BCs:

\begin{itemize}
\item {\bf Dirichlet BC}

\beq \eta(0, t)= D_\rvt(t)
\eeq

\item
{\bf von Neumann BC}
\beq
\partial_x\eta(0,t)= N_\rvt(t)
\eeq
For example, if there were an impermeable wall at $x=0$,

\beq
\partial_x\eta(0,t)=0
\eeq

\item {\bf Mixed (Robin) BC}

\beq
[a \partial_x\eta + b \eta](0, t)=R_\rvt(0,t)
\eeq
\end{itemize}






\section{Notation}
Function $f()$ of one event $(x,t)$
\beq
f(x,t) =f\oneevent{x}{t}
\eeq
Function $G()$ of two events $(x,t)$ and $(x',t')$

\beq
G(x,t;x',t')=G\twoevent{x}{t}{x'}{t'}
\eeq


\beq
\call_{x,t}= \partial_t -D \partial_x^2,\quad
\call G\twoevent{x}{t}{x'}{t'}=\stackt{\delta(x-x')}{\delta(t-t')}
\eeq


\beq
\call^\dagger_{x,t}= -{\partial}_t -D {\partial}_x^2,
\quad G^\dagger\twoevent{x'}{t'}{x}{t}\call^\dagger_{x,t}=\stackt{\delta(x-x')}{\delta(t-t')}
\eeq




\beq
P_\rvx = \int_0^\infty dx'\; \ket{\rvx=x'}\bra{\rvx=x'}
\eeq


\beq
P_\rvt = \int_0^t dt'\; \ket{\rvt=t'}\bra{\rvt=t'}
\eeq

\beq
P_0 = \ket{0}\bra{0}
\eeq

\beq
\partial_{\rvx} =
\int_0^\infty dx'\; \ket{\rvx=x'}\partial_{x'}\bra{\rvx=x'}
\eeq

\beq
\partial_{\rvt} =
\int_0^\infty dt\; \ket{\rvt=t'}\partial_{t'}\bra{\rvt=t'}
\eeq

\beq
\rv{\call}= \partial_\rvt-D\partial^2_\rvx
\eeq

\beq
f(x)\cev{\partial_x} a(x)= a(x) \partial_x f(x)
\eeq
We will use the following suggestive notation
for dual tensor product spaces



\beq \stackt{A}{B}= A\otimes B\eeq

\beq \kett{x}{t}=\ket{x}\otimes\ket{t},
\quad 
\brat{x}{t}=\bra{x}\otimes\bra{t}
\eeq

\section{Master Equation}
\begin{claim}
If $\eta(x,t)$ satisfies
\beq
\call_{x,t}\eta(x,t)=S(x,t)
\eeq
for all $x>0$ and $t>0$,
the initial conditions
\beq
\eta(x,0)=\eta_\rvx(x)
\eeq
 
and has bound

\beq
\ket{\eta}=
G
\stackt{P_\rvx}{P_\rvt}\ket{S}
+G\stackt{P_\rvx}{P_0}\kett{\eta_\rvx}{0}
-D
G\stackt{\left[-\ket{x'}\partial_{x'}\bra{x'}
+
\ket{x'}\cev{\partial}_{x'}\bra{x'}\right]_{x'=0}}{P_\rvt}
\kett{0}{D_\rvt}
\eeq
\end{claim}
\proof

Define $I$, $I_\rvx$ and $I_\rvt$  by
\beq
I = \brat{x}{t}G\rv{\call}
\ket{\eta}-
\brat{x}{t}\rv{\call^\dagger}G^\dagger
\ket{\eta}
\eeq

\beq
I_\rvt =
\brat{x}{t}G
\partial_\rvt
\ket{\eta}-
\brat{x}{t}
(-{\partial}_\rvt)G^\dagger
\ket{\eta}
\eeq

\beq
I_\rvx =
\brat{x}{t}G\partial_\rvx^2
\ket{\eta}-
\brat{x}{t}
\partial_\rvx^2G^\dagger
\ket{\eta}
\eeq
Hence
\beq
\boxed
{I = I_\rvt -D I_\rvx}
\eeq
\hrule
Note that we can simplify the
expression for $I$ as follows:



\beq
\rv{\call}\ket{\eta}=\ket{S},
\quad
\rv{\call^\dagger}G^\dagger=1
\eeq

\beq
\boxed
{I = \brat{x}{t}G\ket{S}-\brat{x}{t}\ket{\eta}
}
\eeq
\hrule
Note that we can simplify the
expression for $I_\rvt$ as follows:


\beqa
I_\rvt &=&
\int_0^t dt'
\brat{x}{t}G
\stackt{1}{\ket{t'}\partial_{t'}\bra{t'}}
\ket{\eta}
+
\brat{x}{t}G
\stackt{1}{\ket{t'}\cev{\partial}_{t'}\bra{t'}}
\ket{\eta}
\eeqa
where we used

\beq
 \bra{t'}G^\dagger\ket{t}=\bra{t}G\ket{t'}
\implies
\partial_{t'} \bra{t'}G^\dagger\ket{t}
=
\bra{t}G\ket{t'}\cev{\partial}_{t'}
\eeq

\beq
\partial_{t'} (\ket{t'}\bra{t'})=
\ket{t'}\partial_{t'}\bra{t'}
+\ket{t'}\cev{\partial}_{t'}\bra{t'}
\eeq


\beqa
I_\rvt &=&
\int_0^t dt'
\brat{x}{t}G
\stackt{1}{\partial_{t'} (\ket{t'}\bra{t'})}
\ket{\eta}
\eeqa

\beq
\int_0^t
dt'\;\partial_{t'} (\ket{t'}\bra{t'})
=
\ket{t}\bra{t}-\ket{0}\bra{0}
\eeq


\beq
\boxed{I_\rvt =
\brat{x}{t}G
\stackt{P_\rvx}{\ket{t}\bra{t}-\ket{0}\bra{0}}
\ket{\eta}}
\eeq
\hrule
Note that we can simplify the
expression for $I_\rvx$ as follows:

\beq
I_\rvx =
\int_0^\infty dx'
\brat{x}{t}G
\stackt{\ket{x'}\partial_{x'}^2\bra{x'}}{1}
\ket{\eta}
-
\brat{x}{t}
G\stackt{\ket{x'}\cev{\partial}_{x'}^2\bra{x'}}{1}
\ket{\eta}
\eeq

\beq
\ket{x'}\partial_{x'}^2\bra{x'}
-\ket{x'}\cev{\partial}_{x'}^2\bra{x'}=
\partial_{x'}\left(\ket{x'}\partial_{x'}\bra{x'}\right)
-\left(\ket{x'}\cev{\partial}_{x'}\bra{x'}\right)\cev{\partial}_{x'}
\eeq

\beq
I_\rvx=
\int_0^\infty dx'
\brat{x}{t}G
\stackt{\partial_{x'}(\ket{x'}\partial_{x'}\bra{x'})
-\left(\ket{x'}\cev{\partial}_{x'}\bra{x'}\right)\cev{\partial}_{x'}}{1}
\ket{\eta}
\eeq

\begin{align}
\int_0^\infty dx'\left[
\partial_{x'}(\ket{x'}\partial_{x'}\bra{x'})
-\left(\ket{x'}\cev{\partial}_{x'}\bra{x'}\right)\cev{\partial}_{x'}\right]
&=
\left[\ket{x'}\partial_{x'}\bra{x'}
-
\ket{x'}\cev{\partial}_{x'}\bra{x'}\right]_0^\infty
\\
&=
-\left[\ket{x'}\partial_{x'}\bra{x'}
-
\ket{x'}\cev{\partial}_{x'}\bra{x'}\right]_{x'=0}
\end{align}

\beq
\boxed{I_\rvx=
\brat{x}{t}G
\stackt{\left[-\ket{x'}\partial_{x'}\bra{x'}
+
\ket{x'}\cev{\partial}_{x'}\bra{x'}\right]_{x'=0}}{1}
\ket{\eta}}
\eeq

\qed
\section{Dirichlet BC}

\begin{claim}\label{cl-de-diriclet-green}
The solution to the equation
\beq
\call \eta (x,t) = S(x,t)\quad
\eeq
for $x>0$ and $t>0$
with initial conditions
\beq
\eta(x,0)= \eta_\rvx(x)
\eeq
and Dirichlet boundary conditions:

\beq
\eta(0,t) = D_\rvt(t)
\eeq
is

\beqa
\eta(x,t)&=&
\int_0^{t} dt'\;
\int_0^\infty dx'\;
G(x,t;x',t')
S(x',t')
\\
&&+
\int_0^{\infty} dx'\;
G(x,t;x', 0)
 \eta_{\rvx}(x')
\\
&&-D
\int_0^{t} dt'\;
[\partial_{x'} G(x,t;x',t')]_{x'=0}
 D_\rvt(t')
\eeqa
where $G(x,t;x',t')$ satisfies
the following equations for $x>0$ and $t>0$


\beq
\call G(x,t;x', t')=\delta(x-x')\delta(t-t')
\eeq
and

\beq
G(x=0,t;x',t')=0
\eeq
\end{claim}
\proof
\qed



Claim \ref{cl-de-diriclet-green} uses a Green function
$G(x,t;x',t')$ with specific properties,
but it does give an explicit expression for $G$
or even prove that such a $G$ exists.
It is easy to check that
the following $G$ has the desired properties

\beq
G = G^- - G^+
\eeq
where

\beq
G^\pm(x,t;x',t')
=
\frac{1}{\sqrt{4\pi D(t-t')}}
\exp\left[-\;\frac{(x\pm x')^2}{4D(t-t')}\right]
\eeq

At the boundary
\beq
\left.
\partial_{x'} G
(x,t;x',t')
\right|_{x'=0}
=
\frac{x}{\sqrt{4\pi D (t-t')^3}}
\exp\left(
-\frac{x^2}{4D(t-t')}
\right)
\eeq
This $G$ was found by the method of images: two
free space solutions $G_\pm$ to $\ul{\call}G_\pm=1$
were superposed so as  to satisfy
$G=0$ at the boundary $x=0$.
\hrule
At first glance, this result appears to imply
that

\beq
\lim_{x\rarrow 0}(-D)
\int_0^{t} dt'\;
[\partial_{x'} G(x,t;x',t')]_{x'=0}
=\delta(t-t')
\eeq
This is indeed the case. Let's prove
explicitly for the Green function $G=G^--G^+$

Suppose

\beq
G^-(x,t;x',t')= P(x-x'; \kappa (t-t'))
\eeq
where $D=\kappa/2$.
Define

\beq
y=\frac{x}{\sqrt{\kappa(t-t')}}=\frac{x}{\ell}
\eeq
Then
\beqa
-D[\partial_{x'} G^-(x,t;x',t')]_{x'=0}&=&
\left(\frac{-\kappa}{2}\right)
\left(\frac{-x}{\ell^2}\right)\frac{1}{\ell\sqrt{2\pi}}\exp\left(
\frac{-x^2}{2\ell^2}\right)
\\
&=&
\frac{x}{2(t-t')}\frac{1}{\ell\sqrt{2\pi}}\exp\left(
\frac{-y^2}{2}\right)
\eeqa

\beq
\ln y = -\frac{1}{2}\ln (t-t')+\cdots\implies \frac{dy}{y}=
-\;\frac{d(t-t')}{2(t-t')}=
\;\frac{dt'}{2(t-t')}
\eeq
Thus,

\beq
dt'\;\frac{x}{2(t-t')}= dy\;\frac{x}{y}=
dy\sqrt{\kappa(t-t')}=dy\ell
\eeq
Hence
\beq
-D\int_0^t dt'[\partial_{x'} G^-(x,t;x',t')]_{x'=0}=
\int_{\frac{x}{\sqrt{\kappa t}}}^\infty
dy \frac{1}{\sqrt{2\pi}}\exp\left(
\frac{-y^2}{2}\right)=\frac{1}{2}
\eeq


It's illuminating to prove this with delta
functions and their derivatives. To do so, assume
\beq
 G^-(x,t;x',t') =\frac{1}{\ell} \delta\left(\frac{x-x'}{\ell}\right),
\eeq
Then

\begin{align}
-D\int_0^t dt'[\partial_{x'} G^-(x,t;x',t')]_{x'=0}
&=
-D\int_0^t \frac{dt'}{\ell^2} \delta'\left(\frac{x}{\ell}\right)
\\
&=
-D\int_0^t\frac{dt'}{\ell^2}\frac{1}{\eps}\left[
\delta\left(\frac{x}{\ell}+\frac{\eps}{2}\right)-
\delta\left(\frac{x}{\ell}-\frac{\eps}{2}\right)
\right]
\\
&=
-
\int_{\frac{x}{\sqrt{\kappa t}}}^\infty \frac{dy}{y}\frac{1}{\eps}\left[
\delta\left(y+\frac{\eps}{2}\right)-
\delta\left(y-\frac{\eps}{2}\right)
\right]
\\
&= \int_{0}^\infty
dy\;\delta\left(y\right)= \frac{1}{2}
\end{align}
Hence, the derivative
of a delta function equals the difference of
2 delta functions. We integrate over half of
one of the two delta functions.


\section{von Neumann BC}

\section{Mixed (Robin) BC}

